\section{Introduction}\label{sec:intro}
% Budget 1.8 pages with Fig 1 (C54).
\begin{figure*}[t]
\centering
% Fig. 1 (hero): the Blender render with the vector overlay (arrows, coordinate lines, scale bars, labels, inset).
% Input from paper/main.tex as % Fig. 1 (hero): the Blender render with the vector overlay (arrows, coordinate lines, scale bars, labels, inset).
% Input from paper/main.tex as % Fig. 1 (hero): the Blender render with the vector overlay (arrows, coordinate lines, scale bars, labels, inset).
% Input from paper/main.tex as \input{../figures/fig1_hero.tex}; the standalone mock defines \figdir, \heroimage and \overlayfile.
% Needs \usepackage{graphicx,tikz} and \usetikzlibrary{arrows.meta}.
\providecommand{\figdir}{../figures/}%
\providecommand{\heroimage}{\figdir fig1_hero.png}%
\providecommand{\overlayfile}{\figdir fig1_overlay.tex}%
\begin{tikzpicture}%
  \node[anchor=south west,inner sep=0pt] at (0,0) {\includegraphics[width=\textwidth]{\heroimage}};%
  \begin{scope}[x=\textwidth,y=\textwidth]%
    \input{\overlayfile}%
  \end{scope}%
\end{tikzpicture}%
; the standalone mock defines \figdir, \heroimage and \overlayfile.
% Needs \usepackage{graphicx,tikz} and \usetikzlibrary{arrows.meta}.
\providecommand{\figdir}{../figures/}%
\providecommand{\heroimage}{\figdir fig1_hero.png}%
\providecommand{\overlayfile}{\figdir fig1_overlay.tex}%
\begin{tikzpicture}%
  \node[anchor=south west,inner sep=0pt] at (0,0) {\includegraphics[width=\textwidth]{\heroimage}};%
  \begin{scope}[x=\textwidth,y=\textwidth]%
    \input{\overlayfile}%
  \end{scope}%
\end{tikzpicture}%
; the standalone mock defines \figdir, \heroimage and \overlayfile.
% Needs \usepackage{graphicx,tikz} and \usetikzlibrary{arrows.meta}.
\providecommand{\figdir}{../figures/}%
\providecommand{\heroimage}{\figdir fig1_hero.png}%
\providecommand{\overlayfile}{\figdir fig1_overlay.tex}%
\begin{tikzpicture}%
  \node[anchor=south west,inner sep=0pt] at (0,0) {\includegraphics[width=\textwidth]{\heroimage}};%
  \begin{scope}[x=\textwidth,y=\textwidth]%
    \input{\overlayfile}%
  \end{scope}%
\end{tikzpicture}%

\IfFileExists{../figures/fig1_caption.tex}{% Caption of Fig. 1 (hero); \input from paper/main.tex and use as \caption{\FigOneCaption}.
\newcommand{\FigOneCaption}{The lone-agent trap at a 0.8~m door (top) and the symmetric pair standoff (bottom), each with its rest-state energy along the coordinate drawn in the scene (right), before and after the drive or tilt that removes the minimum. Agent colour is the stuck-alone margin (cividis, $\pm150$~N); the floor is the backward wall force on a lone 0.25~m agent (log scale, background below 5~N, orange to dark purple up to 1000~N); arrows are the drive (black), the ORCA rest drive (blue), the wall and partner forces (dark magma) and the neighbours' push (grey), on the same force and length scale in both scenes; the vermilion band is the thermal energy $k_BT_{\mathrm{eff}}$; and the wall ranges are $B_w=0.08$~m above and $0.04$~m below.}
}{\newcommand{\FigOneCaption}{\pending{Fig 1 caption}}}
\caption{\FigOneCaption}\label{fig:hero}
\end{figure*}
%
Fatal crowd crushes happen where many people press towards a narrow opening~\cite{liang2024,helbing2012}. Simulations used to plan for such places are only as trustworthy as the steering model that moves each agent, and that includes the reasons the model gives for what it does at a bottleneck.

Steering models are commonly described by the mechanism each one adds: driven particles with repulsive and contact forces in the social force model (SFM)~\cite{helbing1995,helbing2000nature}, a reaction to predicted encounters in the time-to-collision (TTC) force~\cite{karamouzas2014}, and a small linear program over velocity-obstacle constraints in optimal reciprocal collision avoidance (ORCA)~\cite{vandenberg2011}. When a configuration that includes ORCA empties a room that pure SFM leaves jammed, it is tempting to say that reciprocal avoidance resolved the jam. That statement is about mechanism, and a completion rate cannot support it on its own.

Our conference paper~\cite{gang2026} proposed that such statements be tested by pairing a failure-mode ablation with a minimal control: a perturbation of the baseline model that shares none of the paradigm's machinery, tuned until it reproduces the paradigm's headline score. The principle stops short. A matched control shows that a simpler cause suffices for the score, not that the paradigm and the control act in the same way: two physical processes can produce the same completion rate, and one number cannot separate them.

Our thesis is that matching a control to a paradigm's headline completion rate is not evidence of a shared mechanism, and that signatures derived from the physics of the model, measured on the same frozen deadlock states, decide the question. In two deadlock classes of a social-force bottleneck we show that the same headline score hides different mechanisms (Fig.~\ref{fig:hero}). The contributions are:
\begin{itemize}
\item \emph{Method.} A calibrate-then-falsify protocol: a resolution claim is assigned to one of five physically defined mechanism classes, each class predicts how a frozen deadlock responds to a set of probes (Table~\ref{tab:signatures}), and a rest-state lemma removes from consideration every force term that vanishes at rest (Sect.~\ref{sec:signatures}).
\item \emph{Physics.} Narrow-exit deadlocks in the SFM fall into two classes: a lone-agent trap set by the wall force at the door edges, with a stuck-alone criterion that holds in \CriterionAgree\ of \CriterionN\ deadlocks, and a symmetric pair standoff, a strict minimum of the rest-state energy in all \PairStrictMin\ library cases, softest along the common lateral shift, which a lateral bias breaks at a fold. Velocity noise frees the lone trap by collective activated escape over barriers computed from the model's own energy landscape, and contact arches formed at high drive never persist, because the model has no static friction (Sect.~\ref{sec:physics}).
\item \emph{Attribution.} At the lone trap ORCA acts as an added goal drive, deterministic and independent of the time step, whereas noise matched to its rate acts thermally; the two differ in which deadlocks they leave unresolved, in replay and in time-to-free. At the pair class ORCA acts through its constraint part, differently in the two solvers, and a lateral bias matched to its rate frees other standoffs, with a handedness ORCA lacks. In whole runs, noise matched to ORCA's completion rate fails on other seeds, completes later and leaves other final states (Sect.~\ref{sec:attribution}).
\item \emph{Validity.} Rebuilt like for like around the CroMa bottleneck experiments, every model variant passes only \EtenCSigFlowRatio--\EtenCThreeFlowRatio\ of the human flow, with mean longest gaps between passages of \EtenCThreeGap\,s or more against \EtenHumanGap\,s for people (Sect.~\ref{sec:validity}).
\end{itemize}
Sections~\ref{sec:background}--\ref{sec:signatures} give the background, the model as the code implements it, and the classes and signatures; Sects.~\ref{sec:physics}--\ref{sec:validity} the deadlock physics, the attribution results and the comparison with human data.
